% =====================================================================
\chapter{Modulo 7 -- Algoritmi di Ordinamento e Ricerca}
% =====================================================================

% ---------------------------------------------------------------------
\section{Teoria \& Analogia Chiave}
% ---------------------------------------------------------------------

\subsection*{Perché ordinare}

Ordinare un array significa disporre i suoi elementi in ordine crescente (o
decrescente). Sembra un problema banale, ma è uno dei più studiati
dell'informatica, per due motivi:
\begin{itemize}
  \item i dati ordinati si \textbf{cercano} enormemente più in fretta: la ricerca
        binaria su un milione di elementi ordinati richiede 20 confronti, quella
        lineare un milione (Esercizio~6.9);
  \item molti altri problemi diventano facili sui dati ordinati: trovare i
        duplicati, il mediano, gli elementi comuni a due elenchi (Esercizio~7.10).
\end{itemize}
In questo modulo vediamo cinque algoritmi di ordinamento: tre semplici e lenti
($O(n^2)$) e due ``intelligenti'' e veloci ($O(n \log n)$), basati sul divide et
impera del Modulo 6. Per ognuno devi sapere: \textbf{l'idea}, \textbf{il codice}
e \textbf{la complessità}.

Due proprietà che compaiono nei quiz:
\begin{itemize}
  \item un algoritmo è \textbf{stabile} se due elementi con la stessa chiave
        restano nell'ordine in cui erano. Esempio: ordino per voto degli studenti
        già in ordine alfabetico; con un algoritmo stabile, a parità di voto restano
        in ordine alfabetico (Esercizio~7.6);
  \item un algoritmo è \textbf{in loco} (\emph{in place}) se usa solo una quantità
        costante di memoria in più oltre all'array.
\end{itemize}

\subsection*{Bubble sort (ordinamento a bolle)}

\textbf{Idea.} Scorri l'array confrontando gli elementi \emph{vicini}, \cc|a[j]| e
\cc|a[j+1]|: se sono nell'ordine sbagliato li scambi. Alla fine della prima passata
l'elemento più grande è ``salito'' fino in fondo, come una bolla nell'acqua. Ripeti:
la seconda passata porta in fondo il secondo più grande, e così via. Dopo $n-1$
passate l'array è ordinato.

\begin{lstlisting}
for (int i = 0; i < n - 1; i++) {            // n-1 passate
    bool scambiato = false;
    for (int j = 0; j < n - 1 - i; j++)      // gli ultimi i sono già a posto
        if (a[j] > a[j + 1]) { scambia(a[j], a[j + 1]); scambiato = true; }
    if (!scambiato) break;                   // nessuno scambio: già ordinato
}
\end{lstlisting}

\textbf{Complessità.} Due cicli annidati ``a triangolo'': circa $n^2/2$ confronti,
quindi $O(n^2)$ nel caso peggiore (array al contrario) e medio. Con il controllo
\cc|scambiato|, se l'array è già ordinato basta una passata: caso migliore $O(n)$.
Stabile e in loco.

\subsection*{Selection sort (ordinamento per selezione)}

\textbf{Idea.} Cerca il \textbf{minimo} di tutto l'array e scambialo con il primo
elemento. Poi cerca il minimo tra gli elementi dal secondo in poi e mettilo in
seconda posizione. E così via: alla passata $i$ selezioni il minimo della parte
non ancora ordinata e lo metti in posizione $i$.

\begin{lstlisting}
for (int i = 0; i < n - 1; i++) {
    int iMin = i;
    for (int j = i + 1; j < n; j++)
        if (a[j] < a[iMin]) iMin = j;
    scambia(a[i], a[iMin]);
}
\end{lstlisting}

\textbf{Complessità.} I confronti sono sempre $n(n-1)/2$, \emph{qualunque} sia
l'input: $O(n^2)$ in tutti i casi, anche se l'array è già ordinato. Fa però
pochissimi scambi (al più $n-1$). In loco, non stabile nella versione con scambi.

\subsection*{Insertion sort (ordinamento per inserimento)}

\begin{analogia}[: le carte in mano]
È il modo in cui ordini le carte da gioco: le prendi una alla volta e
\textbf{inserisci} ognuna al posto giusto tra quelle che hai già in mano, facendo
scorrere a destra quelle più grandi.
\end{analogia}

\begin{lstlisting}
for (int i = 1; i < n; i++) {
    int x = a[i];                  // la carta da inserire
    int j = i - 1;
    while (j >= 0 && a[j] > x) {   // sposto a destra le carte più grandi
        a[j + 1] = a[j];
        j--;
    }
    a[j + 1] = x;                  // la metto nel buco rimasto
}
\end{lstlisting}

\textbf{Complessità.} Caso peggiore (array al contrario): ogni elemento attraversa
tutti i precedenti, $O(n^2)$. Caso migliore (già ordinato): il \cc|while| non parte
mai, $O(n)$. Su array \emph{quasi} ordinati è molto veloce. Stabile e in loco. Tra i
tre algoritmi semplici è il migliore nella pratica (nel laboratorio fa circa metà
dei confronti del bubble sort).

\subsection*{Merge sort (ordinamento per fusione)}

È il primo algoritmo \textbf{divide et impera}, ed è \textbf{ricorsivo} (la
ricorsione si approfondisce nel Modulo 8).
\begin{passopasso}[: merge sort]
\begin{enumerate}
  \item \textbf{Dividi}: taglia l'array a metà.
  \item \textbf{Conquista}: ordina ricorsivamente la metà sinistra e la metà destra.
        Il caso base è un pezzo con 0 o 1 elementi, che è già ordinato.
  \item \textbf{Combina}: \textbf{fondi} (\emph{merge}) le due metà ordinate in un
        unico array ordinato.
\end{enumerate}
\end{passopasso}

Il cuore è la \textbf{fusione}. Hai due mazzetti di carte già ordinati, a faccia in
su: guardi le due carte in cima, prendi la più piccola e la metti nel mazzo
risultato. Ripeti finché uno dei due mazzetti finisce, poi accodi il resto
dell'altro. Ogni carta viene toccata una volta: fondere costa $O(n)$. Serve però un
\textbf{array di appoggio} dove scrivere il risultato: il merge sort \textbf{non} è
in loco.

\begin{lstlisting}
void merge(int a[], int sx, int centro, int dx) {
    int* b = new int[dx - sx + 1];          // appoggio
    int i = sx, j = centro + 1, k = 0;
    while (i <= centro && j <= dx)          // finché entrambe le metà hanno elementi
        b[k++] = (a[i] <= a[j]) ? a[i++] : a[j++];
    while (i <= centro) b[k++] = a[i++];    // avanzi della sinistra
    while (j <= dx) b[k++] = a[j++];        // avanzi della destra
    for (k = 0; k < dx - sx + 1; k++) a[sx + k] = b[k];
    delete[] b;
}
void mergeSort(int a[], int sx, int dx) {
    if (sx >= dx) return;                   // caso base
    int centro = (sx + dx) / 2;
    mergeSort(a, sx, centro);
    mergeSort(a, centro + 1, dx);
    merge(a, sx, centro, dx);
}
\end{lstlisting}

\textbf{Complessità.} È la ricorrenza $C_N = 2C_{N/2} + N$ del Modulo 6: $\log n$
livelli di dimezzamento, lavoro $n$ per livello, quindi $O(n \log n)$ \textbf{in
tutti i casi} (migliore, medio e peggiore), perché divide sempre a metà
indipendentemente dai dati. Spazio $O(n)$ per l'appoggio. È \textbf{stabile} grazie
al \cc|<=| nella fusione: a parità prende prima l'elemento di sinistra.

\subsection*{Quick sort (ordinamento rapido)}

\begin{passopasso}[: quick sort]
\begin{enumerate}
  \item Scegli un elemento, il \textbf{pivot} (perno). Nella versione delle slide è
        l'elemento centrale.
  \item \textbf{Partiziona}: sposta a sinistra tutti gli elementi minori del pivot e a
        destra tutti quelli maggiori. Ora ogni elemento della parte sinistra è
        $\le$ di ogni elemento della parte destra.
  \item Ordina ricorsivamente la parte sinistra e la parte destra. Non serve nessuna
        fusione: le due parti sono già ``in ordine tra loro''.
\end{enumerate}
\end{passopasso}

La partizione delle slide usa due indici che si vengono incontro: \cc|i| parte da
sinistra e avanza finché trova un elemento che \emph{non} è minore del pivot (sta
dalla parte sbagliata); \cc|j| parte da destra e arretra finché trova un elemento che
non è maggiore del pivot. Quando entrambi si fermano, scambi \cc|a[i]| e \cc|a[j]|
e prosegui. Quando \cc|i| supera \cc|j| la partizione è finita.

\begin{lstlisting}
void quickSort(int a[], int sx, int dx) {
    int i = sx, j = dx;
    int pivot = a[(sx + dx) / 2];
    while (i <= j) {
        while (a[i] < pivot) i++;
        while (a[j] > pivot) j--;
        if (i <= j) { scambia(a[i], a[j]); i++; j--; }
    }
    if (sx < j) quickSort(a, sx, j);    // ci sono ancora elementi a sinistra
    if (i < dx) quickSort(a, i, dx);    // ci sono ancora elementi a destra
}
\end{lstlisting}

\textbf{Complessità.} Dipende da quanto bene il pivot divide l'array:
\begin{itemize}
  \item \textbf{caso migliore e medio}: il pivot divide in due parti più o meno
        uguali, quindi $\log n$ livelli da $n$: $O(n \log n)$. Le slide precisano
        che basta una divisione ``proporzionale'' (anche 10\%--90\%) perché resti
        $O(n \log n)$;
  \item \textbf{caso peggiore}: il pivot è sempre il minimo o il massimo, le parti
        hanno $0$ e $n-1$ elementi e i livelli diventano $n$: $O(n^2)$. Succede ad
        esempio scegliendo come pivot il \emph{primo} elemento di un array già
        ordinato (Esercizio~7.7: 1\,999\,000 confronti contro 20\,010).
\end{itemize}
In pratica è il più veloce di tutti (nel laboratorio batte il merge sort in tempo),
perché lavora in loco e fa operazioni semplici. Non è stabile.

\subsection*{Riepilogo da memorizzare}

\begin{center}\small
\begin{tabular}{lllllcc}
\toprule
\textbf{Algoritmo} & \textbf{Migliore} & \textbf{Medio} & \textbf{Peggiore} & \textbf{Spazio extra} & \textbf{Stabile} & \textbf{In loco} \\
\midrule
Bubble    & $O(n)$       & $O(n^2)$      & $O(n^2)$      & $O(1)$      & sì & sì \\
Selection & $O(n^2)$     & $O(n^2)$      & $O(n^2)$      & $O(1)$      & no & sì \\
Insertion & $O(n)$       & $O(n^2)$      & $O(n^2)$      & $O(1)$      & sì & sì \\
Merge     & $O(n\log n)$ & $O(n\log n)$  & $O(n\log n)$  & $O(n)$      & sì & no \\
Quick     & $O(n\log n)$ & $O(n\log n)$  & $O(n^2)$      & $O(\log n)$ (stack) & no & sì \\
\bottomrule
\end{tabular}
\end{center}
Ricorda dal Modulo 6 che \textbf{nessun} algoritmo che ordina confrontando gli
elementi può fare meglio di $\Omega(n \log n)$ nel caso peggiore: merge sort è
quindi \textbf{ottimo}. Si può fare meglio solo \emph{senza} confronti, quando si
sa qualcosa sui valori: ad esempio il counting sort per valori interi in un
intervallo piccolo, come i voti da 18 a 30 (Esercizio~7.8).

\subsection*{Ricerca lineare e ricerca binaria}

\textbf{Ricerca lineare}: controlla gli elementi uno a uno. Funziona su qualunque
array, ordinato o no. Costo $O(n)$ nel caso peggiore e medio.

\textbf{Ricerca binaria} (o dicotomica): funziona \textbf{solo su array
ordinati}. Guarda l'elemento centrale: se è quello cercato hai finito; se è più
piccolo, il valore può stare solo nella metà destra; se è più grande, solo nella metà
sinistra. Butti via metà dell'array e ripeti. È la rubrica del Modulo 6: $O(\log n)$.

\begin{lstlisting}
int ricercaBinaria(const int a[], int n, int x) {
    int sx = 0, dx = n - 1;
    while (sx <= dx) {                   // finché l'intervallo non è vuoto
        int m = (sx + dx) / 2;
        if (a[m] == x) return m;
        if (a[m] < x) sx = m + 1;        // butto la metà sinistra (m compreso)
        else dx = m - 1;                 // butto la metà destra (m compreso)
    }
    return -1;                           // non trovato
}
\end{lstlisting}
Se devi cercare \emph{molte} volte nello stesso array, conviene ordinarlo una volta
($O(n \log n)$) e poi fare tante ricerche binarie ($O(\log n)$ ciascuna).

\begin{inparole}
\begin{itemize}
  \item Bubble: scambi tra vicini, il massimo ``sale'' in fondo. Selection: scegli il
        minimo e mettilo davanti. Insertion: inserisci come le carte.
        Tutti $O(n^2)$.
  \item Merge: dividi a metà, ordina le metà, fondi. Sempre $O(n \log n)$, serve un
        array d'appoggio, stabile.
  \item Quick: pivot, partiziona, ricorri. $O(n \log n)$ in media, $O(n^2)$ con pivot
        sfortunati, in loco, il più veloce in pratica.
  \item Ricerca binaria: solo su array ordinati, dimezza a ogni passo, $O(\log n)$.
\end{itemize}
\end{inparole}

% ---------------------------------------------------------------------
\section{Rappresentazione Visiva della Memoria (TikZ)}
% ---------------------------------------------------------------------

\textbf{Bubble sort: la prima passata.} In giallo la coppia confrontata; il
massimo (82) arriva in fondo.
\begin{center}
\begin{tikzpicture}
  \foreach \riga/\vals/\hl in {0/{38,27,43,3,9,82,10}/0, 1/{27,38,43,3,9,82,10}/2, 2/{27,38,3,43,9,82,10}/3, 3/{27,38,3,9,43,82,10}/4, 4/{27,38,3,9,43,10,82}/9} {
    \foreach \v [count=\i from 0] in \vals {
      \pgfmathsetmacro{\evid}{(\i==\hl || \i==\hl+1) ? 1 : 0}
      \ifdim\evid pt>0.5pt
        \node[cellina,fill=yellow!40!sfondo] at (\i*0.8,-\riga*0.7) {\v};
      \else
        \node[cellina,fill=sfondo] at (\i*0.8,-\riga*0.7) {\v};
      \fi
    }
  }
  \node[etichetta,anchor=west,text width=5cm] at (6,-1.4) {si confrontano coppie vicine e si scambiano se fuori ordine; righe intermedie omesse};
  \node[cellina,fill=green!25!sfondo] at (4.8,-2.8) {82};
\end{tikzpicture}
\end{center}

\needspace{7cm}
\textbf{Merge sort sull'esempio delle slide.} In alto le divisioni, in basso le
fusioni.
\begin{center}
\begin{tikzpicture}[blocco/.style={draw,fill=stackcol,font=\ttfamily\footnotesize,minimum height=0.55cm,inner sep=2pt}]
  \node[blocco] (a) at (0,0) {38 27 43 3 9 82 10};
  \node[blocco] (b1) at (-2.4,-0.9) {38 27 43 3};
  \node[blocco] (b2) at (2.4,-0.9) {9 82 10};
  \node[blocco] (c1) at (-3.4,-1.8) {38 27};
  \node[blocco] (c2) at (-1.4,-1.8) {43 3};
  \node[blocco] (c3) at (1.6,-1.8) {9 82};
  \node[blocco] (c4) at (3.4,-1.8) {10};
  \foreach \x/\v [count=\i] in {-3.9/38,-2.9/27,-1.9/43,-0.9/3,1.1/9,2.1/82}
    \node[blocco] (d\i) at (\x,-2.7) {\v};
  \draw (a) -- (b1); \draw (a) -- (b2);
  \draw (b1) -- (c1); \draw (b1) -- (c2); \draw (b2) -- (c3); \draw (b2) -- (c4);
  \draw (c1) -- (d1); \draw (c1) -- (d2); \draw (c2) -- (d3); \draw (c2) -- (d4); \draw (c3) -- (d5); \draw (c3) -- (d6);
  \node[blocco,fill=heapcol] (e1) at (-3.4,-3.7) {27 38};
  \node[blocco,fill=heapcol] (e2) at (-1.4,-3.7) {3 43};
  \node[blocco,fill=heapcol] (e3) at (1.6,-3.7) {9 82};
  \node[blocco,fill=heapcol] (e4) at (3.4,-3.7) {10};
  \node[blocco,fill=heapcol] (f1) at (-2.4,-4.6) {3 27 38 43};
  \node[blocco,fill=heapcol] (f2) at (2.4,-4.6) {9 10 82};
  \node[blocco,fill=green!20!sfondo] (g) at (0,-5.5) {3 9 10 27 38 43 82};
  \draw (d1) -- (e1); \draw (d2) -- (e1); \draw (d3) -- (e2); \draw (d4) -- (e2); \draw (d5) -- (e3); \draw (d6) -- (e3); \draw (c4) -- (e4);
  \draw (e1) -- (f1); \draw (e2) -- (f1); \draw (e3) -- (f2); \draw (e4) -- (f2);
  \draw (f1) -- (g); \draw (f2) -- (g);
  \node[etichetta,anchor=west] at (4.6,-1.3) {DIVIDI};
  \node[etichetta,anchor=west] at (4.6,-4.6) {FONDI (merge)};
\end{tikzpicture}
\end{center}

\needspace{5cm}
\textbf{Quick sort: una partizione.} Pivot = elemento centrale (3). \cc|i| avanza
finché \cc|a[i] < 3|, \cc|j| arretra finché \cc|a[j] > 3|; poi si scambia.
\begin{center}
\begin{tikzpicture}
  \foreach \v [count=\i from 0] in {38,27,43,3,9,82,10}
    \node[cellina,fill=sfondo] (q\i) at (\i*0.8,0) {\v};
  \node[cellina,fill=yellow!40!sfondo] at (q3) {3};   % il pivot evidenziato
  \draw[freccia] ($(q0.south)+(0,-0.5)$) node[below,font=\small] {i} -- (q0.south);
  \draw[freccia] ($(q3.south)+(0,-0.5)$) node[below,font=\small] {j} -- (q3.south);
  \node[etichetta,anchor=west,text width=7.6cm] at (6,0) {
    \texttt{a[0] = 38} non è $< 3$: \texttt{i} si ferma subito. \texttt{j} arretra da 10
    fino a \texttt{a[3] = 3}, che non è $> 3$. Si scambiano: \texttt{3 27 43 38 9 82 10}.
    Alla fine a sinistra c'è solo \texttt{3}, a destra tutto il resto: si ricorre sulla destra.};
\end{tikzpicture}
\end{center}

% ---------------------------------------------------------------------
\section{Pattern di Codice di Laboratorio}
% ---------------------------------------------------------------------

I cinque algoritmi e la ricerca binaria, scritti nella forma più compatta e
leggibile, provati sullo stesso array delle slide (con un 27 doppio, per verificare
i duplicati). I tre algoritmi semplici hanno la stessa firma e sono chiamati tramite
un array di puntatori a funzione.

\codice{codice/m07/pattern.cpp}
\uscita{codice/m07/output/pattern.txt}

% ---------------------------------------------------------------------
\section{Tips \& Trappole da Esame}
% ---------------------------------------------------------------------

\begin{trappola}[: ricerca binaria su array non ordinato]
Compila e gira, ma dà risultati sbagliati senza nessun errore. La ricerca binaria
richiede che l'array sia \textbf{ordinato}: se non lo è, ordinalo prima.
\end{trappola}

\begin{trappola}[: indici fuori di uno]
I bug più comuni: \cc|j < n - i| invece di \cc|j < n - 1 - i| nel bubble sort (si
legge \cc|a[n]|), \cc|while (sx < dx)| invece di \cc|<=| nella ricerca binaria (salta
l'ultimo elemento), \cc|mergeSort(a, 0, n)| invece di \cc|n - 1|. Prova sempre con
array di 0, 1 e 2 elementi.
\end{trappola}

\begin{trappola}[: l'array d'appoggio del merge]
Un array locale \cc|int b[n];| con \cc|n| variabile non è C++ standard (alcuni
compilatori lo accettano) e con $n$ grande esaurisce lo stack. Usa \cc|new int[...]|
e \cc|delete[]|, oppure allocalo una volta sola e passalo alle chiamate (come nel
laboratorio).
\end{trappola}

\begin{trappola}[: complessità ``tutti i casi'']
Domande tipiche: il merge sort è $O(n \log n)$ \textbf{anche} nel caso peggiore; il
quick sort no ($O(n^2)$). Il selection sort è $O(n^2)$ \textbf{anche} nel caso
migliore; bubble (con il flag) e insertion no ($O(n)$).
\end{trappola}

\begin{consiglio}[: overload con l'involucro]
Merge sort e quick sort lavorano su \cc|(a, sx, dx)|. Scrivi anche una versione
\cc|ordina(int a[], int n)| che chiama \cc|mergeSort(a, 0, n - 1)|: così chi usa la
funzione non deve preoccuparsi degli indici.
\end{consiglio}

\begin{consiglio}[: l'operatore virgola]
Nel laboratorio \cc|while (j >= 0 && (confronti++, a[j] > x))| usa l'operatore
virgola: esegue \cc|confronti++| e poi vale \cc|a[j] > x|. Serve solo a contare i
confronti senza cambiare l'algoritmo; nel codice d'esame non serve.
\end{consiglio}

% ---------------------------------------------------------------------
\section{Quiz a Crocette (Prova Scritta)}
% ---------------------------------------------------------------------

\begin{quiz}{Qual è la complessità del merge sort nel caso peggiore?}
  \item $O(n)$
  \item $O(n \log n)$
  \item $O(n^2)$
  \item $O(\log n)$
\end{quiz}

\begin{quiz}{Qual è la complessità del quick sort nel caso peggiore?}
  \item $O(n \log n)$
  \item $O(n)$
  \item $O(n^2)$
  \item $O(2^n)$
\end{quiz}

\begin{quiz}{Quale algoritmo esegue sempre lo stesso numero di confronti, anche se l'array è già ordinato?}
  \item bubble sort con controllo sugli scambi
  \item insertion sort
  \item selection sort
  \item nessuno
\end{quiz}

\begin{quiz}{Cosa richiede la ricerca binaria per funzionare correttamente?}
  \item un array di numeri interi
  \item un array ordinato
  \item un array di dimensione pari
  \item un array senza duplicati
\end{quiz}

\begin{quiz}{Quale affermazione sul merge sort è \textbf{falsa}?}
  \item segue il paradigma divide et impera
  \item è stabile
  \item ordina in loco senza memoria aggiuntiva
  \item la fusione di due metà ordinate costa tempo lineare
\end{quiz}

\begin{quiz}{Nel quick sort, cosa garantisce la fase di partizione?}
  \item che le due parti abbiano la stessa dimensione
  \item che ogni elemento della parte sinistra sia $\le$ di ogni elemento della parte destra
  \item che la parte sinistra sia già ordinata
  \item che il pivot sia il mediano
\end{quiz}

\begin{quiz}{Un algoritmo di ordinamento si dice stabile se:}
  \item non va mai in crash
  \item ha la stessa complessità in tutti i casi
  \item mantiene l'ordine relativo degli elementi con chiave uguale
  \item non usa memoria aggiuntiva
\end{quiz}

\begin{quiz}{Quanti confronti servono al più per trovare un elemento con la ricerca binaria in un array ordinato di 1000 elementi?}
  \item circa 10
  \item circa 100
  \item circa 500
  \item 1000
\end{quiz}

\begin{quiz}{Quale algoritmo ha complessità $O(n)$ nel caso migliore (array già ordinato)?}
  \item merge sort
  \item selection sort
  \item insertion sort
  \item quick sort
\end{quiz}

\begin{quiz}{Si può ordinare per confronti con complessità $O(n)$ nel caso peggiore?}
  \item sì, con il quick sort
  \item sì, con il bubble sort ottimizzato
  \item no: ogni ordinamento per confronti è $\Omega(n \log n)$ nel caso peggiore
  \item no, il limite è $\Omega(n^2)$
\end{quiz}

\begin{soluzione}[: risposte commentate]
\begin{enumerate}[label=\textbf{\arabic*.}]
  \item \textbf{B}. Divide sempre a metà: $O(n \log n)$ in ogni caso.
  \item \textbf{C}. Con pivot sempre estremo le partizioni sono sbilanciate.
  \item \textbf{C}. Cerca comunque il minimo in tutta la parte restante.
  \item \textbf{B}. I duplicati (D) non sono un problema.
  \item \textbf{C}. Serve un array d'appoggio di dimensione $n$.
  \item \textbf{B}. Le dimensioni possono essere molto diverse (A falsa); le due
        parti vanno poi ordinate ricorsivamente (C falsa).
  \item \textbf{C}.
  \item \textbf{A}. $2^{10} = 1024$: circa 10 dimezzamenti.
  \item \textbf{C}. Il \cc|while| interno non parte mai: un confronto per elemento.
        Vale anche per il bubble sort con il flag, che però non è tra le risposte.
  \item \textbf{C}. Limite inferiore del problema (Modulo 6). Il counting sort è
        lineare perché \emph{non} confronta.
\end{enumerate}
\end{soluzione}

% ---------------------------------------------------------------------
\section{Laboratorio a Step (Simulazione d'Esame)}
% ---------------------------------------------------------------------

\begin{traccia}[: prestazioni degli algoritmi di ordinamento (dalle slide)]
Implementare un metodo \cc|generaArrayCasuale(int* a, int num)| e gli algoritmi
di ordinamento; generare array di dimensioni crescenti per potenze di 10 e
confrontare le prestazioni degli algoritmi al crescere della dimensione, misurando
il tempo con \cc|<chrono>| come indicato nelle slide.
\end{traccia}

\begin{step}{gli strumenti}
Scrivi \cc|generaArrayCasuale|, \cc|copia| (per ordinare ogni volta la
\emph{stessa} sequenza) e \cc|bool ordinato(const int* a, int n)| per verificare il
risultato. Dichiara un contatore globale di confronti.\\[0.3em]
\textbf{Concetto nuovo:} misurare in modo equo (stessi dati per tutti).
\end{step}

\begin{step}{due algoritmi quadratici}
Implementa bubble sort (con uscita anticipata) e insertion sort, incrementando il
contatore a ogni confronto tra elementi.\\[0.3em]
\textbf{Concetto nuovo:} algoritmi $O(n^2)$.
\end{step}

\begin{step}{merge sort}
Implementa il merge sort allocando l'array d'appoggio \textbf{una sola volta} in una
funzione involucro \cc|mergeSort(int a[], int n)| e passandolo alle chiamate
ricorsive.\\[0.3em]
\textbf{Concetto nuovo:} divide et impera, fusione, memoria aggiuntiva $O(n)$.
\end{step}

\begin{step}{quick sort}
Implementa il quick sort con pivot centrale e partizione a due indici, come nelle
slide, più l'involucro \cc|quickSort(int a[], int n)|.\\[0.3em]
\textbf{Concetto nuovo:} partizione, ordinamento in loco.
\end{step}

\begin{step}{il confronto}
Scrivi \cc|misura(nome, puntatoreAFunzione, array, n)|, che copia l'array, azzera
il contatore, cronometra l'ordinamento con \cc|high_resolution_clock| e stampa
confronti e millisecondi. Usala per $n$ = $10^2$, $10^3$, $10^4$, $10^5$ (i quadratici
solo fino a $10^4$).\\[0.3em]
\textbf{Concetto nuovo:} misurare i tempi reali; puntatori a funzione come parametri.
\end{step}

\subsection*{Soluzione completa}
\codice{codice/m07/lab.cpp}
\uscita{codice/m07/output/lab.txt}

\begin{consiglio}[: cosa notare nella soluzione]
\begin{itemize}
  \item I \textbf{tempi} cambiano a ogni esecuzione e da PC a PC; i
        \textbf{confronti} no. Per questo il contatore è più affidabile per studiare la
        complessità.
  \item Da $n = 1000$ a $n = 10\,000$ il bubble sort passa da circa 0{,}5 a 50 milioni
        di confronti ($\times 100$): quadratico. Merge e quick crescono di poco più di
        $\times 10$: $n \log n$.
  \item Il quick sort fa \emph{più} confronti del merge sort ma impiega \emph{meno}
        tempo: non copia avanti e indietro in un array d'appoggio.
  \item \cc|setw(10)| (da \cc|<iomanip>|) stampa il numero in una colonna larga 10
        caratteri, così la tabella resta allineata.
\end{itemize}
\end{consiglio}

% ---------------------------------------------------------------------
\section{Domande e Risposte Modello per l'Orale}
% ---------------------------------------------------------------------

\begin{orale}[D1. Come funziona il merge sort e qual è la sua complessità?]
È un algoritmo divide et impera: divide l'array in due metà, le ordina
ricorsivamente fino ad arrivare a sottoarray di un elemento (caso base) e poi le
fonde con la procedura merge, che scorre le due metà ordinate e prende ogni volta il
minore dei due elementi in testa. La fusione costa $O(n)$ e i livelli di ricorsione
sono $\log n$, perciò la complessità è $O(n \log n)$ nel caso migliore, medio e
peggiore. Richiede un array d'appoggio, quindi spazio $O(n)$, ed è stabile.
\end{orale}

\begin{orale}[D2. Come funziona il quick sort? Qual è il suo caso peggiore?]
Si sceglie un pivot e si partiziona l'array in modo che a sinistra ci siano gli
elementi minori o uguali al pivot e a destra quelli maggiori o uguali; poi si
ordinano ricorsivamente le due parti, senza bisogno di fusione. Se le partizioni sono
bilanciate (anche solo in proporzione costante) la complessità è $O(n \log n)$, che è
anche il caso medio. Il caso peggiore si ha quando il pivot è sempre il minimo o il
massimo: le partizioni hanno dimensione $0$ e $n-1$ e la complessità diventa
$O(n^2)$, ad esempio con pivot sul primo elemento di un array già ordinato. Ordina in
loco e in pratica è il più veloce.
\end{orale}

\begin{orale}[D3. Differenza tra bubble, selection e insertion sort?]
Sono tutti $O(n^2)$ nel caso peggiore. Il bubble sort scambia elementi adiacenti
fuori ordine, portando a ogni passata il massimo in fondo; con il controllo sugli
scambi termina in $O(n)$ se l'array è ordinato. Il selection sort cerca a ogni
passata il minimo della parte non ordinata e lo mette in posizione: fa sempre
$n(n-1)/2$ confronti ma pochi scambi. L'insertion sort inserisce ogni elemento al suo
posto tra i precedenti già ordinati: $O(n)$ nel caso migliore ed efficiente su array
quasi ordinati. Bubble e insertion sono stabili.
\end{orale}

\begin{orale}[D4. Come funziona la ricerca binaria e perché è logaritmica?]
Si applica a un array ordinato. Si confronta il valore cercato con l'elemento
centrale: se coincide si è trovato; se è maggiore si prosegue nella metà destra,
altrimenti nella sinistra. A ogni passo l'intervallo di ricerca si dimezza, quindi
dopo al più circa $\log_2 n$ confronti l'intervallo diventa vuoto. La ricorrenza è
$C_N = C_{N/2} + 1$, con soluzione $O(\log n)$.
\end{orale}

\begin{orale}[D5. Cosa significa che un ordinamento è stabile? E in loco?]
È stabile se elementi con la stessa chiave mantengono, dopo l'ordinamento, l'ordine
relativo che avevano prima: ad esempio studenti già in ordine alfabetico, ordinati
per voto, restano in ordine alfabetico a parità di voto. Merge, insertion e bubble
sort sono stabili; selection e quick sort no. È in loco se usa solo memoria
aggiuntiva costante (o al più lo stack della ricorsione): tutti tranne il merge sort,
che ha bisogno di un array d'appoggio.
\end{orale}

% ---------------------------------------------------------------------
\section{Esercizi al PC}
% ---------------------------------------------------------------------

\begin{esercizio}{Bubble sort e passate}{facile}
Modifica il bubble sort con uscita anticipata in modo che restituisca il numero di
passate eseguite. Provalo su un array già ordinato, su uno quasi ordinato (due
elementi scambiati) e su uno al contrario.
\end{esercizio}

\begin{esercizio}{Selection sort decrescente}{facile}
Scrivi un selection sort che ordina in modo \textbf{decrescente} (cerca il massimo)
e stampa l'array dopo ogni passata.
\end{esercizio}

\begin{esercizio}{Insertion sort di stringhe}{facile}
Ordina in ordine alfabetico un array di \cc|string| con l'insertion sort. Inserisci
anche un nome che inizia con la minuscola: dove finisce, e perché?
\end{esercizio}

\begin{esercizio}{Ricerca binaria al rallentatore}{facile}
Scrivi la ricerca binaria iterativa stampando a ogni passo \cc|sx|, \cc|dx|, \cc|m| e
\cc|a[m]|. Cerca un valore presente e uno assente.
\end{esercizio}

\begin{esercizio}{Ricerca binaria ricorsiva}{medio}
Scrivi la ricerca binaria in versione ricorsiva e una funzione
\cc|posizioneInserimento(a, n, x)| che restituisce la prima posizione in cui
\cc|x| si può inserire mantenendo l'ordine (con duplicati: la prima occorrenza).
\end{esercizio}

\begin{esercizio}{Merge sort stabile di studenti}{medio}
Ordina per voto con il merge sort un array di \cc|struct Studente| (nome, voto) già
in ordine alfabetico, e verifica che a parità di voto l'ordine alfabetico resti.
Cosa succederebbe usando \cc|<| invece di \cc|<=| nella fusione?
\end{esercizio}

\begin{esercizio}{Il caso peggiore del quick sort}{medio}
Implementa un quick sort con pivot sul \textbf{primo} elemento e confronta i suoi
confronti con quelli della versione delle slide (pivot centrale) su un array già
ordinato di 2000 elementi.
\end{esercizio}

\begin{esercizio}{Counting sort dei voti}{medio}
Ordina 20 voti casuali tra 18 e 30 senza confronti: conta quante volte compare ogni
voto in un array di contatori, poi riscrivi l'array. Che complessità ha?
\end{esercizio}

\begin{esercizio}{Ordinare con un criterio a scelta}{difficile}
Scrivi \cc|template <typename T> void ordina(T a[], int n, bool (*minore)(const T&, const T&))|
e usalo per ordinare un array di \cc|Persona| per età, per nome e per età
decrescente, e un array di \cc|int| in ordine decrescente, passando ogni volta una
funzione di confronto diversa.
\end{esercizio}

\begin{esercizio}{Sfruttare l'ordine}{difficile}
Su array ordinati scrivi in $O(n)$: (a) una funzione che trova due elementi con somma
data usando due indici che si vengono incontro; (b) una funzione che trova gli elementi
comuni a due array, con un procedimento simile alla fusione.
\end{esercizio}

\subsection*{Soluzioni}
Le soluzioni commentate, con il codice e l'output reale, sono nel PDF separato
\textbf{Soluzioni} (\texttt{soluzioni-light.pdf} / \texttt{soluzioni-dark.pdf}),
allo stesso numero di esercizio. Il codice è anche in \texttt{codice/m07/}.
